\chapter{2013年上海卷（秋文）}

\section{填空题}

\begin{problem}
不等式 \(\frac{x}{2x - 1} < 0\) 的解为\fillinblank{}.

\begin{answer}
\(0 < x < \frac{1}{2}\)
\end{answer}
\end{problem}

\begin{problem}
在等差数列 \(\{a_n\}\) 中，若 \(a_1 + a_2 + a_3 + a_4 = 30\)，则 \(a_2 + a_3 =\fillinblank{}.\)

\begin{answer}
\(15\)
\end{answer}
\end{problem}

\begin{problem}
设 \(m \in \R\)，\(m^2 + m - 2 + \left(m^2 - 1\right)\i\) 是纯虚数，其中 \(\i\) 是虚数单位，则 \(m =\fillinblank{}.\)

\begin{answer}
\(-2\)
\end{answer}
\end{problem}

\begin{problem}
若 \(\begin{vmatrix} x & 2 \\ 1 & 1 \end{vmatrix} = 0\)，\(\begin{vmatrix} x & y \\ 1 & 1 \end{vmatrix} = 1\)，则 \(x + y =\fillinblank{}.\)

\begin{answer}
\(3\)
\end{answer}

\begin{solution}
由第一个行列式等于零，得 \(x-2=0\)，所以 \(x=2\). 由第二个行列式等于 \(1\)，得 \(x-y=1\)，代入 \(x=2\) 得 \(y=1\). 因此 \(x+y=2+1=3\).
\end{solution}
\end{problem}

\begin{problem}
已知 \(\triangle ABC\) 的内角 \(A\)，\(B\)，\(C\) 所对的边分别是 \(a\)，\(b\)，\(c\). 若 \(a^2 + ab + b^2 - c^2 = 0\)，则角 \(C\) 的大小是\fillinblank{}.

\begin{answer}
\(\frac{2\pi}{3}\)
\end{answer}
\end{problem}

\begin{problem}
某学校高一年级男生人数占该年级学生人数的 \(40\%\). 在一次考试中，男、女生平均分数分别为 \(75\)，\(80\)，则这次考试该年级学生平均分数为\fillinblank{}.

\begin{answer}
\(78\)
\end{answer}
\end{problem}

\begin{problem}
设常数 \(a \in \R\). 若 \(\left(x^2 + \frac{a}{x}\right)^5\) 的二项展开式中 \(x^7\) 项的系数为 \(-10\)，则 \(a =\fillinblank{}.\)

\begin{answer}
\(-2\)
\end{answer}
\end{problem}

\begin{problem}
方程 \(\frac{9}{3^x - 1} + 1 = 3^x\) 的实数解为\fillinblank{}.

\begin{answer}
\(x = \log_{3}4\)
\end{answer}
\end{problem}

\begin{problem}
若 \(\cos x\cos y+\sin x\sin y=\frac{1}{3}\)，则 \(\cos\left(2x-2y\right)=\)\fillinblank{}.

\begin{answer}
\(-\frac{7}{9}\)
\end{answer}
\end{problem}

\begin{problem}
已知圆柱 \(\Omega\) 的母线长为 \(l\)，底面半径为 \(r\)，\(O\) 是上底面圆心，\(A\)，\(B\) 是下底面圆周上两个不同的点，\(BC\) 是母线，如图. 若直线 \(OA\) 与 \(BC\) 所成角的大小为 \(\frac{\pi}{6}\)，则 \(\frac{l}{r} =\fillinblank{}.\)

\bitmapfigure[max width=0.36\linewidth]{img/2013/shanghai_liberal/q10_fig1.png}

\begin{answer}
\(\sqrt{3}\)
\end{answer}
\end{problem}

\begin{problem}
盒子中装有编号为 \(1\)，\(2\)，\(3\)，\(4\)，\(5\)，\(6\)，\(7\) 的七个球，从中任意取出两个，则这两个球的编号之积为偶数的概率是\fillinblank{}.（结果用最简分数表示）

\begin{answer}
\(\frac{6}{7}\)
\end{answer}
\end{problem}

\begin{problem}
设 \(AB\) 是椭圆 \(\Gamma\) 的长轴，点 \(C\) 在 \(\Gamma\) 上，且 \(\angle CBA = \frac{\pi}{4}\)，若 \(AB = 4\)，\(BC = \sqrt{2}\)，则 \(\Gamma\) 的两个焦点之间的距离为\fillinblank{}.

\begin{answer}
\(\frac{4\sqrt{6}}{3}\)
\end{answer}
\end{problem}

\begin{problem}
设常数 \(a > 0\). 若 \(9x + \frac{a^2}{x} \geqslant a + 1\) 对一切正实数 \(x\) 成立，则 \(a\) 的取值范围为\fillinblank{}.

\begin{answer}
\(a \geqslant \frac{1}{5}\)
\end{answer}
\end{problem}

\begin{problem}
已知正方形 \(ABCD\) 的边长为 \(1\). 记以 \(A\) 为起点，其余顶点为终点的向量分别为 \(\bs{a}_1\)，\(\bs{a}_2\)，\(\bs{a}_3\)；以 \(C\) 为起点，其余顶点为终点的向量分别为 \(\bs{c}_1\)，\(\bs{c}_2\)，\(\bs{c}_3\). 若 \(i\)，\(j\)，\(k\)，\(l \in \{1, 2, 3\}\) 且 \(i \neq j\)，\(k \neq l\)，则 \(\left(\bs{a}_i + \bs{a}_j\right) \cdot \left(\bs{c}_k + \bs{c}_l\right)\) 的最小值是\fillinblank{}.

\begin{answer}
\(-5\)
\end{answer}
\end{problem}

\section{选择题}

\begin{problem}
函数 \(f(x) = x^2 - 1\) （\(x \geqslant 1\)）的反函数为 \(f^{-1}(x)\)，则 \(f^{-1}(2)\) 的值是（\(\quad\)）

\choices
    {\(\sqrt{3}\)}
    {\(-\sqrt{3}\)}
    {\(1 + \sqrt{2}\)}
    {\(1 - \sqrt{2}\)}
\begin{answer}
A
\end{answer}
\end{problem}

\begin{problem}
设常数 \(a \in \R\)，集合 \(A = \{x \mid \left(x - 1\right)\left(x - a\right) \geqslant 0\}\)，\(B = \{x \mid x \geqslant a - 1\}\). 若 \(A \cup B = \R\)，则 \(a\) 的取值范围为（\(\quad\)）

\choices
    {\(\left(-\infty, 2\right)\)}
    {\(\left(-\infty, 2\right]\)}
    {\(\left(2, +\infty\right)\)}
    {\(\left[2, +\infty\right)\)}
\begin{answer}
B
\end{answer}
\end{problem}

\begin{problem}
钱大姐常说“便宜没好货”，她这句话的意思是：“不便宜”是“好货”的（\(\quad\)）

\choices
    {充分条件}
    {必要条件}
    {充分必要条件}
    {既非充分也非必要条件}
\begin{answer}
B
\end{answer}
\end{problem}

\begin{problem}
记椭圆 \(\frac{x^2}{4} + \frac{ny^2}{4n + 1} = 1\) 围成的区域（含边界）为 \(\Omega_n\) （\(n = 1, 2, \cdots\)），当点 \(\left(x, y\right)\) 分别在 \(\Omega_1\)，\(\Omega_2\)，\(\cdots\) 上时，\(x + y\) 的最大值分别是 \(M_1\)，\(M_2\)，\(\cdots\)，则 \(\displaystyle\lim_{n \to \infty} M_n =\)（\(\quad\)）

\choices
    {\(0\)}
    {\(\frac{1}{4}\)}
    {\(2\)}
    {\(2\sqrt{2}\)}
\begin{answer}
D
\end{answer}
\end{problem}

\section{解答题}

\begin{problem}
如图，正三棱锥 \(O - ABC\) 底面边长为 \(2\)，高为 \(1\)，求该三棱锥的体积及表面积.

\bitmapfigure[max width=0.42\linewidth]{img/2013/shanghai_liberal/q19_fig1.png}

\begin{solution}
正三角形 \(ABC\) 的高为 \(\sqrt{2^2-1^2}=\sqrt{3}\)，所以底面积为 \(\frac{1}{2}\times2\times\sqrt{3}=\sqrt{3}\). 因此三棱锥的体积为
\[
V=\frac{1}{3}\times\sqrt{3}\times1=\frac{\sqrt{3}}{3}.
\]
设 \(H\) 是正三角形 \(ABC\) 的中心，\(D\) 是 \(AB\) 的中点. 则 \(HD\) 是正三角形的内切半径，\(HD=\frac{\sqrt{3}}{3}\)，且正三棱锥的高 \(OH=1\)，所以侧面三角形 \(OAB\) 的高为
\[
OD=\sqrt{OH^2+HD^2}=\sqrt{1^2+\left(\frac{\sqrt{3}}{3}\right)^2}=\frac{2\sqrt{3}}{3}.
\]
每个侧面三角形的面积为 \(\frac{1}{2}\times2\times\frac{2\sqrt{3}}{3}=\frac{2\sqrt{3}}{3}\)，三个侧面的总面积为 \(2\sqrt{3}\). 故三棱锥的表面积为
\[
S=\sqrt{3}+2\sqrt{3}=3\sqrt{3}.
\]
\end{solution}
\end{problem}

\begin{problem}
甲厂以 \(x\text{ 千克/小时}\) 的速度匀速生产某种产品（生产条件要求 \(1 \leqslant x \leqslant 10\)），每一小时可获得的利润是 \(100\left(5x + 1 - \frac{3}{x}\right)\text{ 元}\).

\begin{enumerate}
\item 求证：生产 \(a\text{ 千克}\) 该产品所获得的利润为 \(100a\left(5 + \frac{1}{x} - \frac{3}{x^2}\right)\text{ 元}\)；
\item 要使生产 \(900\text{ 千克}\) 该产品获得的利润最大，问：甲厂应该选取何种生产速度？并求最大利润.
\end{enumerate}

\begin{solution}
\begin{enumerate}
\item 生产 \(a\text{ 千克}\) 产品需要 \(\frac{a}{x}\) 小时，所以所获得的利润为
\[
\frac{a}{x}\times100\left(5x+1-\frac{3}{x}\right)=100a\left(5+\frac{1}{x}-\frac{3}{x^2}\right)\text{ 元}.
\]
\item 当生产 \(900\text{ 千克}\) 产品时，利润函数为
\[
P(x)=90000\left(5+\frac{1}{x}-\frac{3}{x^2}\right),\quad 1\leqslant x\leqslant10.
\]
其导数为
\[
P'(x)=90000\left(-\frac{1}{x^2}+\frac{6}{x^3}\right)=\frac{90000(6-x)}{x^3}.
\]
由于 \(x>0\)，当 \(1\leqslant x<6\) 时 \(P'(x)>0\)，当 \(6<x\leqslant10\) 时 \(P'(x)<0\)，所以 \(P(x)\) 在 \([1,6]\) 上递增，在 \([6,10]\) 上递减，最大值在 \(x=6\) 时取得. 最大利润为
\[
P(6)=90000\left(5+\frac{1}{6}-\frac{3}{36}\right)=90000\times\frac{61}{12}=457500\text{ 元}.
\]
因此甲厂应选择 \(6\text{ 千克/小时}\) 的生产速度，最大利润为 \(457500\text{ 元}\).
\end{enumerate}
\end{solution}
\end{problem}

\begin{problem}
已知函数 \(f(x) = 2\sin \omega x\). 其中常数 \(\omega > 0\).

\begin{enumerate}
\item 令 \(\omega = 1\)，判断函数 \(F(x) = f(x) + f\left(x + \frac{\pi}{2}\right)\) 的奇偶性并说明理由；
\item 令 \(\omega = 2\)，将函数 \(y = f(x)\) 的图象向左平移 \(\frac{\pi}{6}\) 个单位，再向上平移 \(1\) 个单位，得到函数 \(y = g(x)\) 的图象. 对任意的 \(a \in \R\)，求 \(y = g(x)\) 在区间 \([a, a + 10\pi]\) 上零点个数的所有可能值.
\end{enumerate}

\begin{solution}
\begin{enumerate}
\item 当 \(\omega=1\) 时，\(f(x)=2\sin x\)，所以
\[
F(x)=2\sin x+2\sin\left(x+\frac{\pi}{2}\right)=2\sin x+2\cos x.
\]
于是 \(F(-x)=2\cos x-2\sin x\). 又 \(F(0)=2\neq-2=-F(0)\)，故 \(F(x)\) 不是奇函数；并且 \(F\left(\frac{\pi}{2}\right)=2\)，\(F\left(-\frac{\pi}{2}\right)=-2\)，故 \(F(x)\) 不是偶函数. 因此 \(F(x)\) 既不是奇函数，也不是偶函数.
\item 当 \(\omega=2\) 时，\(f(x)=2\sin2x\). 图象向左平移 \(\frac{\pi}{6}\) 后，横坐标 \(x\) 处的函数值为 \(f\left(x+\frac{\pi}{6}\right)\)，再向上平移 \(1\) 个单位，故
\[
g(x)=2\sin\left(2x+\frac{\pi}{3}\right)+1.
\]
令 \(g(x)=0\)，得
\[
\sin\left(2x+\frac{\pi}{3}\right)=-\frac{1}{2}.
\]
因此
\[
2x+\frac{\pi}{3}=-\frac{\pi}{6}+2k\pi\quad\text{或}\quad2x+\frac{\pi}{3}=\frac{7\pi}{6}+2k\pi,\quad k\in\Z,
\]
即零点为
\[
x=-\frac{\pi}{4}+k\pi\quad\text{或}\quad x=\frac{5\pi}{12}+k\pi,\quad k\in\Z.
\]
函数 \(g(x)\) 的周期为 \(\pi\)，每个周期内有两个零点. 区间 \([a,a+10\pi]\) 的长度是 \(10\) 个周期，因此当 \(a\) 不是零点时，区间内有 \(20\) 个零点；当 \(a\) 是零点时，\(a+10\pi\) 也是零点，两个端点同时计入，区间内有 \(21\) 个零点. 故所有可能值为 \(20\)，\(21\).
\end{enumerate}
\end{solution}
\end{problem}

\begin{problem}
已知函数 \(f(x) = 2 - |x|\)，无穷数列 \(\{a_n\}\) 满足 \(a_{n+1} = f(a_n)\)，\(n \in \N^*\).

\begin{enumerate}
\item 若 \(a_1 = 0\)，求 \(a_2\)，\(a_3\)，\(a_4\)；
\item 若 \(a_1 > 0\)，且 \(a_1\)，\(a_2\)，\(a_3\) 成等比数列，求 \(a_1\) 的值；
\item 是否存在 \(a_1\)，使得 \(a_1\)，\(a_2\)，\(a_3\)，\(\cdots\)，\(a_n\)，\(\cdots\) 成等差数列？若存在，求出所有这样的 \(a_1\)；若不存在，说明理由.
\end{enumerate}

\begin{solution}
\begin{enumerate}
\item 当 \(a_1=0\) 时，\(a_2=2-|0|=2\)，\(a_3=2-|2|=0\)，\(a_4=2-|0|=2\).
\item 设 \(a_1=t>0\).

当 \(0<t\leqslant2\) 时，\(a_2=2-t\geqslant0\)，从而 \(a_3=2-(2-t)=t\). 由 \(a_1,a_2,a_3\) 成等比数列，得 \((2-t)^2=t^2\)，
所以 \(4-4t=0\)，即 \(t=1\).

当 \(t>2\) 时，\(a_2=2-t<0\)，从而 \(a_3=2-|2-t|=4-t\). 于是 \((2-t)^2=t(4-t)\)，
即 \(2t^2-8t+4=0\)，解得 \(t=2\pm\sqrt{2}\). 结合 \(t>2\)，得 \(t=2+\sqrt{2}\). 因此 \(a_1=1\) 或 \(a_1=2+\sqrt{2}\).
\item 假设数列成等差数列，公差为 \(d\)，则 \(a_n=a_1+(n-1)d\). 由于 \(a_{n+1}=2-|a_n|\leqslant2\)，若 \(d>0\)，则 \(a_n\) 当 \(n\) 足够大时大于 \(2\)，矛盾. 若 \(d<0\)，则 \(a_n\) 当 \(n\) 足够大时为负数，此时递推式变为 \(a_{n+1}=2+a_n\)，所以 \(d=2\)，又与 \(d<0\) 矛盾. 因此只能有 \(d=0\).

此时所有 \(a_n\) 都等于 \(a_1\)，由递推式得 \(a_1=2-|a_1|\). 当 \(a_1\geqslant0\) 时，\(a_1=2-a_1\)，得 \(a_1=1\)；当 \(a_1<0\) 时，\(a_1=2+a_1\)，矛盾. 因此存在且唯一的取值为 \(a_1=1\).
\end{enumerate}
\end{solution}
\end{problem}

\begin{problem}
如图，已知双曲线 \(C_1: \frac{x^2}{2} - y^2 = 1\)，曲线 \(C_2: |y| = |x| + 1\). \(P\) 是平面内一点，若存在过点 \(P\) 的直线 \(l\) 与 \(C_1\)，\(C_2\) 都有公共点，则称 \(P\) 为“\(C_1 - C_2\) 型点”.

\begin{enumerate}
\item 在正确证明 \(C_1\) 的左焦点是“\(C_1 - C_2\) 型点”时，要使用一条过该焦点的直线，试写出一条这样的直线的方程（不要求验证）；
\item 设直线 \(y = kx\) 与 \(C_2\) 有公共点. 求证 \(|k| > 1\)，进而证明原点不是“\(C_1 - C_2\) 型点”；
\item 求证：圆 \(x^2 + y^2 = \frac{1}{2}\) 内的点都不是“\(C_1 - C_2\) 型点”.
\end{enumerate}

\FigureLayoutDeclare{img/2013/shanghai_liberal/q23_fig1.png}{6.5cm}{0bp 15bp 92bp 0bp}
\bitmapfigure[max width=0.38\linewidth]{img/2013/shanghai_liberal/q23_fig1.png}

\begin{solution}
\begin{enumerate}
\item 双曲线 \(C_1\) 的左焦点为 \(F=(-\sqrt{3},0)\)，直线 \(y=\frac{x}{\sqrt{3}}+1\) 经过点 \(F\)，且经过 \(C_2\) 上的点 \((0,1)\)，可作为所需直线.
\item 设直线 \(y=kx\) 与 \(C_2\) 的公共点为 \((x,y)\). 若 \(x=0\)，则由 \(y=kx\) 得 \(y=0\)，这与 \(C_2\) 的方程 \(|y|=|x|+1\) 矛盾，所以 \(x\neq0\). 于是
\[
|k|=\frac{|y|}{|x|}=\frac{|x|+1}{|x|}>1.
\]
若原点是 \(C_1-C_2\) 型点，则存在过原点的直线同时与 \(C_1\)，\(C_2\) 有公共点. 直线 \(x=0\) 与 \(C_1\) 联立得 \(-y^2=1\)，无实数解；其他过原点的直线可写为 \(y=kx\). 由上面结论，这样的直线若与 \(C_2\) 相交则 \(|k|>1\)，而它与 \(C_1\) 相交需满足
\[
\frac{x^2}{2}-k^2x^2=1,
\]
从而 \(\left(\frac{1}{2}-k^2\right)x^2=1\)，必有 \(k^2<\frac{1}{2}\)，即 \(|k|<\frac{1}{\sqrt{2}}\)，矛盾. 因此原点不是 \(C_1-C_2\) 型点.
\item 设点 \(P=(x_0,y_0)\) 在圆 \(x^2+y^2=\frac{1}{2}\) 的内部，即 \(x_0^2+y_0^2<\frac{1}{2}\). 反设 \(P\) 是 \(C_1-C_2\) 型点，设过 \(P\) 且与两条曲线都有公共点的直线为 \(\ell\).

若 \(\ell\) 为竖直直线 \(x=c\)，由于它与 \(C_1\) 相交，必有 \(c^2/2\geqslant1\)，即 \(|c|\geqslant\sqrt{2}>\frac{1}{\sqrt{2}}\). 所以 \(\ell\) 上任一点到原点的距离不可能小于 \(\frac{1}{\sqrt{2}}\)，与 \(P\) 在圆内矛盾.

若 \(\ell\) 非竖直，设其方程为 \(y=kx+b\). 当 \(|k|\leqslant1\) 时，取 \(\ell\) 与 \(C_2\) 的公共点 \((x,y)\)，则
\[
|b|=|y-kx|\geqslant\bigl||y|-|k||x|\bigr|=1+(1-|k|)|x|\geqslant1.
\]
因此原点到 \(\ell\) 的距离 \(d\) 满足
\[
d=\frac{|b|}{\sqrt{1+k^2}}\geqslant\frac{1}{\sqrt{1+k^2}}\geqslant\frac{1}{\sqrt{2}}.
\]

当 \(|k|>1\) 时，将 \(y=kx+b\) 代入 \(C_1\)，得
\[
\left(\frac{1}{2}-k^2\right)x^2-2kbx-(b^2+1)=0.
\]
该方程有实数解，所以其判别式满足
\[
\Delta=2\left(b^2+1-2k^2\right)\geqslant0,
\]
从而 \(b^2\geqslant2k^2-1\). 于是
\[
d^2=\frac{b^2}{1+k^2}\geqslant\frac{2k^2-1}{1+k^2}\geqslant\frac{1}{2},
\]
其中最后一个不等号等价于 \(k^2\geqslant1\). 综上，任何同时与 \(C_1\)，\(C_2\) 相交的直线到原点的距离都不小于 \(\frac{1}{\sqrt{2}}\). 但 \(P\in\ell\)，所以该距离不大于 \(OP=\sqrt{x_0^2+y_0^2}<\frac{1}{\sqrt{2}}\)，矛盾. 因此圆内的点都不是 \(C_1-C_2\) 型点.
\end{enumerate}
\end{solution}
\end{problem}
